% ==============================================================================
% ssec:cs_ksvs_bankovnictvi --- Případová studie: KSVS bankovnictví a pojišťovnictví
% Zdroje: KSVS_OSPPP_2026, CMKOS_ZpravaKV2025,
%         zakon_zp_2006, zakon_kv_1991, eu_directive_2022_2041
% ==============================================================================

\Acf{KSVS} pro odvětví bankovnictví a~pojišťovnictví na~období 2026--2027 byla podepsána 16.~prosince 2025 v~Praze mezi \acf{SBP} a~\acf{OSPPP}. Smlouva se~vztahuje na~zaměstnavatele v~odvětvích \acgen{NACE}~64, 65 a~66 sdružené v~\acloc{SBP}. Vzhledem k~tomu, že u~této \acgen{KSVS} nedošlo k~extenzi závaznosti podle \S~7 \acgen{ZKV}, je~smlouva právně závazná výhradně pro zaměstnavatele sdružené v~signatářské organizaci -- podle statistik \acgen{SBP} jde o~16~velkých zaměstnavatelů z~řad bank a~pojišťoven s~více než \SI{46000}{\person}. Zajišťuje tak pokrytí významné části odvětvové pracovní síly v~\acloc{ČR}. Odvětví představuje protipól maloobchodu: vysoká přidaná hodnota finančních služeb a~kvalifikace pracovníků motivují obě strany k~investicím do~stabilizačních nadstandardů.\par

Minimální základní hrubá měsíční mzda činí \SI{28000}{\czk\per\month} po~odpracování šesti měsíců nepřetržitého \aclgen{PP} u~zaměstnavatele. To je o~\SI{25}{\percent} nad zákonnou minimální mzdou pro rok 2026 (\SI{22400}{\czk\per\month}). Výjimky z~této garantované částky se ovšem vztahují na~nekvalifikované manuální pozice, trainee programy, přepážkové pracovníky pojišťoven a~zaměstnance odměňované provizemi. Mzdové příplatky překračují zákonná minima, přičemž nejvýraznějším nadstandardem je příplatek za~práci v~sobotu a~v~neděli ve~výši \SI{50}{\percent} \aclgen{PV} (pětinásobek zákonem předepsaného minima), který s~ohledem na~výjimečnost víkendové práce ve~finančním sektoru plní zejména funkci drahé kompenzace pohotovostních zásahů.\par

Zásadní inovaci představuje čl.~7 upravující práci z~domova. Smlouva striktně vyžaduje oboustrannou dobrovolnost a~písemnou dohodu podle \S~317 \aclgen{ZP}. Zaměstnavateli ukládá povinnost zajistit veškerou pracovní techniku a~zabezpečenou komunikaci, hradit paušální náhradu nákladů na~energie a~internet nejméně v~rozsahu limitu \S~190a \aclgen{ZP} a~vynucuje pravidlo výhradně osobního výkonu práce ze~strany zaměstnance kvůli ochraně osobních údajů a~klientskému tajemství. Tato úprava reaguje na specifikum odvětví, kde podíl home office podle \acgen{ČSÚ} přesahuje \SI{30}{\percent} zaměstnanců a~je nejvyšší v~rámci národního trhu práce.\par

Sociální fond je sjednán jako obligatorní s~tvorbou nejméně \SI{2,5}{\percent} z~ročních mzdových prostředků, doplněný o~příspěvek na~životní pojištění ve~výši \SI{0,5}{\percent}. Fond slouží ke~krytí nadlimitního stravování, rekreačních a~kulturních potřeb a~sociální výpomoci. Sektorové benefity dále zahrnují tři dny placeného zdravotního volna (sick days) ročně nahrazené podle~základní mzdy a~jeden den placeného volna za~čtvrtletí pro~těhotné zaměstnankyně a~pro osamělé rodiče s~dětmi mladšími patnácti let, což odráží vysoký podíl žen v~odvětví (podíl žen dosahuje přibližně \SI{60}{\percent}).\par

Pracovní podmínky jsou nahlíženy i~skrze psychosociální rizika a~vysokou kognitivní zátěž v~důsledku tlaku na~výkon, digitalizace a~odpovědnosti za~citlivé údaje. Smlouva proto zavazuje zaměstnavatele k~předcházení šikany a~podpoře profesního vzdělávání.\par

Celková bilance této \acgen{KSVS} vykazuje nejsilnější postavení zaměstnanců mezi analyzovanými případy. Dvouleté smluvní období poskytuje zaměstnavateli plánovací stabilitu a~nadstandardní benefity účinně chrání kvalifikovanou pracovní sílu. Vysoká přidaná hodnota finančního sektoru a~silná organizovanost jak na straně \acgen{SBP}, tak na straně \acgen{OSPPP} umožňují dosáhnout partnerského vyjednávání s~vysokým standardem krytí.~\cite{KSVS_OSPPP_2026}~\cite{CMKOS_ZpravaKV2025}\par
